\documentclass[11pt,a4paper]{article}
\usepackage[utf8]{inputenc}
\usepackage[russian,english]{babel}
\usepackage{amsmath,amssymb,amsthm}
\usepackage{mathtools}
\usepackage{geometry}
\usepackage{booktabs}
\usepackage{graphicx}
\usepackage{hyperref}
\usepackage{longtable}
\usepackage{array}
\usepackage{caption}
\geometry{margin=2.5cm}

\newtheorem{definition}{Definition}
\newtheorem{proposition}{Proposition}
\newtheorem{theorem}{Theorem}
\newtheorem{remark}{Remark}

\title{Continuous Topological Accessibility (RFA) as a 3D-Aware GRN Filter: \\
Anticorrelation with eQTL and No Prediction of $\Delta$RNA}
\author{Timofeev Andrey \\
\small AI Center for SCO+ Countries, Saint Petersburg, Russia}
\date{October 2026}

\begin{document}
\maketitle

\begin{abstract}
We formalize \emph{Relative Forest Accessibility} (RFA), a topological accessibility index computed as the solution of $(I - \gamma P)x = v_0$, where $P$ is a Markov transition matrix built from Hi-C contacts and $v_0$ is a vector of TF peaks (or ATAC peaks). RFA yields a continuous accessibility profile across all genomic bins.

We test RFA on three targets. First, on promoter TF counts (GM12878 chr1, 10\,kb, 811 genes), RFA loses to 1D distance (AUC 0.836 vs 0.976); this is an artifact because the target and the 1D distance look at the same window. Second, on eQTL (GTEx v10 LCL, chr1+17+19, 10\,kb, 1625 genes), RFA \emph{anticorrelates} with eQTL: $\rho(\mathrm{RFA}, -\log_{10} p_{\min}) = -0.09$, $\mathrm{AUC}(\mathrm{RFA}) = 0.386$ (below chance), $\mathrm{AUC}(-\mathrm{RFA}) = 0.614$. High RFA corresponds to fewer eQTL, and dose-response analysis confirms that $\mathrm{AUC}(\mathrm{RFA})$ decreases with eQTL strength (mean trend $-0.056$, 19/20 TFs). Third, on $\Delta$RNA (MCF10A vs MCF7, 40\,kb, ATAC peaks), RFA does not predict $\Delta$RNA ($\rho = +0.005$, permutation $p=0.87$, MDR$=0.067$ at $n=1730$); within-line sanity checks are also zero.

We conclude that RFA anticorrelates with eQTL and does not predict $\Delta$RNA. The anticorrelation may reflect expression buffering: genes with strong eQTL are less accessible in 3D. RFA is not a universal predictor but a filter with an inverted sign on eQTL and no signal on $\Delta$RNA.
\end{abstract}

\section{Introduction}

\subsection{Problem}
Existing 3D-aware TF$\rightarrow$gene filters (SSD 2017, Esearch3D 2023, ChIPSP 2026) use either a binary (loop yes/no) or a scalar (direct contact, 1D distance). RFA offers a continuous accessibility profile. The key question is: \emph{where does RFA give signal, and in which direction}.

\subsection{Positioning}
We position RFA not as a ``new framework'' but as ``continuous topological accessibility as a resolution-dependent 3D-aware GRN filter''. RFA is not a universal predictor but a filter with a clearly delineated applicability domain and a \emph{signal direction}.

\subsection{Contributions}
\begin{enumerate}
    \item Formalization of RFA via a Markov random walk (Section~\ref{sec:theory}).
    \item \textbf{Anticorrelation of RFA with eQTL} ($\mathrm{AUC}(-\mathrm{RFA})=0.614$, dose-response trend $-0.056$).
    \item \textbf{Negative result on $\Delta$RNA} ($\rho=0.005$, MDR$=0.067$).
    \item \textbf{Methodological warning:} an automatic sign flip in AUC can mask anticorrelation as ``high AUC''; dose-response reveals it.
\end{enumerate}

\section{Theory}\label{sec:theory}

We formalize the computation of a \emph{topological accessibility index} (RFA) that quantifies how accessible a genomic locus is to a transcription factor (TF) given the 3D chromatin architecture measured by Hi-C. The construction proceeds in four stages: from Hi-C contact counts to a weighted graph; from the graph to a Markov transition matrix; from the transition matrix to a continuous accessibility profile; from the profile to per-gene scores and validation statistics. Throughout, the Markov chain is a \emph{mathematical instrument for deconvolving spatial noise and accounting for multi-step paths}, not a kinetic model of TF motion.

\subsection{Hi-C Graph Construction}

Let $C \in \mathbb{R}_{\geq 0}^{n \times n}$ be a symmetric Hi-C contact matrix at resolution $r$ (e.g., $r = 10\,\mathrm{kb}$), where $n$ is the number of genomic bins  measured by Hi-C \cite{dekker2013,rao2014}. We apply:
\begin{enumerate}
    \item \textbf{Diagonal zeroing}: $C_{ii} := 0$ for all $i$, eliminating self-contacts that would otherwise dominate the degree distribution.
    \item \textbf{Top-$k$ sparsification}: For each row $i$, retain only the $k$ largest entries, yielding a sparse matrix $\tilde{A}$. This replaces dense noise with the strongest topological signals. Statistical treatment of Hi-C noise and sparsification is discussed
in \cite{hu2013}.
    \item \textbf{Symmetrization}: $A := \max(\tilde{A}, \tilde{A}^\top)$, preserving symmetry.
\end{enumerate}

A/B compartments (the dominant large-scale chromatin structure) are captured by the first principal component of $\log(1+C)$. To isolate fine-scale topological structure from compartment signal, we regress each row of $A$ on the compartment score and retain residuals:
\begin{equation}
A_{ij}^{\mathrm{res}} = A_{ij} - \big(\beta_0^{(i)} + \beta_1^{(i)} \mathrm{PC1}_j\big), \qquad A_{ij}^{\mathrm{res}} := \max(A_{ij}^{\mathrm{res}}, 0),
\end{equation}
where $(\beta_0^{(i)}, \beta_1^{(i)})$ are obtained by ordinary least squares on the support of row $i$.

\subsection{Markov Transition Matrix}

Define the degree vector $d_i = \sum_j A_{ij}$ and the diagonal matrix $D = \mathrm{diag}(d_1, \dots, d_n)$. The row-stochastic transition matrix is
\begin{equation}
P = D^{-1} A, \qquad P_{ij} = \frac{A_{ij}}{d_i}, \qquad \sum_j P_{ij} = 1.
\end{equation}
The matrix $P$ encodes the probability that a random walker at bin $i$ moves to bin $j$ in one step.

\begin{remark}
The Markov chain is \emph{not} a kinetic model of TF motion. It is a graph-theoretic object whose powers $P^t$ aggregate the contributions of all paths of length $t$ through the chromatin graph.
\end{remark}

\begin{proposition}
The matrix $P$ has the following properties:
\begin{enumerate}
    \item $P$ is row-stochastic: $P \mathbf{1} = \mathbf{1}$.
    \item The stationary distribution $\pi$ satisfies $\pi P = \pi$, with $\pi_i = d_i / \sum_j d_j$ (degree-weighted).
    \item $P$ is similar to a symmetric matrix: $D^{1/2} P D^{-1/2} = D^{-1/2} A D^{-1/2}$, which is symmetric. Hence all eigenvalues of $P$ are real and lie in $[-1, 1]$.
\end{enumerate}
\end{proposition}

\subsection{Relative Forest Accessibility (RFA)}

Given a source distribution $v_0 \in \mathbb{R}_{\geq 0}^n$ (e.g., TF binding peaks in bins), we define the RFA profile as
\begin{equation}
x_\gamma := (I - \gamma P)^{-1} v_0 = \sum_{t=0}^{\infty} \gamma^t P^t v_0,
\label{eq:rfa}
\end{equation}
where $\gamma \in (0,1)$ is a damping factor.

\begin{proposition}[Convergence]
The Neumann series \eqref{eq:rfa} converges for all $\gamma \in (0,1)$.
\end{proposition}
\begin{proof}
Since $\|P\|_\infty = 1$ and $\rho(P) \leq 1$, we have $\|(\gamma P)^t\|_\infty \leq \gamma^t$, and $\sum_{t=0}^\infty \gamma^t = 1/(1-\gamma) < \infty$.
\end{proof}

The profile $x_\gamma$ has three equivalent interpretations:
\begin{enumerate}
    \item \textbf{Path sum}: $x_{\gamma,i} = \sum_{t=0}^\infty \gamma^t (P^t v_0)_i$ is a weighted sum over all paths of length $t$ from peaks to bin $i$, with exponential decay in path length.
    \item \textbf{Steady-state with restart}: $x_\gamma$ is the unique solution of $(I - \gamma P) x = v_0$, equivalent to a random walk with restart (RWR) at rate $1-\gamma$.
    \item \textbf{Forest accessibility}: $x_{\gamma,i}$ is proportional to the total weight of spanning forests of the graph rooted at the peak set, with edge weights weighted by $\gamma$.
\end{enumerate}

The parameter $\gamma$ controls the \emph{effective path length}. For $\gamma \to 0$, $x_\gamma \to v_0$ (peaks only). For $\gamma \to 1$, $x_\gamma$ approaches the stationary distribution $\pi$ (global mixing). We tested $\gamma \in \{0.5, 0.7, 0.85, 0.95\}$; at 10\,kb with TF peaks, $\gamma = 0.85$ works best. At 40\,kb with ATAC peaks, RFA gives no signal (MDR$=0.067$), consistent with the hypothesis that the signal lives at fine scales and is sensitive to the type of $v_0$.

\subsection{Per-Gene Scoring and Baselines}

Let $\mathcal{G}$ be the set of genes with a promoter annotation. For each gene $g$, let $b(g)$ be the bin containing its promoter ($b(g) = \lfloor \mathrm{start}(g)/r \rfloor$). Define
\begin{equation}
\mathrm{RFA}(g) := x_{\gamma, b(g)}.
\end{equation}

We compare RFA against three baselines:
\begin{enumerate}
    \item \textbf{1D distance}: $D_{1D}(g) = \min_{p \in \mathrm{peaks}} |\mathrm{start}(g) - p|$.
    \item \textbf{Direct contact (ChIPSP proxy)}: $\mathrm{DC}(g) = \max_{p \in \mathrm{peaks}} A_{b(g), p}$.
    \item \textbf{Combined model}: $\mathrm{Comb}(g) = z(\mathrm{RFA}(g)) + z(\mathrm{DC}(g))$, where $z(\cdot)$ is z-score normalization.
\end{enumerate}

For a continuous target $y_g$, we report Spearman $\rho(\mathrm{RFA}, y)$ and partial $\rho(\mathrm{RFA}, y \mid D_{1D}, \mathrm{DC})$, obtained by rank-residualizing both variables against the controls. For binary classification $y_g \in \{0,1\}$, we report ROC-AUC, with DeLong-style bootstrap ($B=1000$) for $\Delta \mathrm{AUC}$.

\subsection{Statistical Properties}

\textbf{Why does RFA outperform direct contact on promoter count?} At high resolution (10\,kb), direct contacts are noisy: each bin contains few counts, and $\mathrm{DC}(g)$ has high variance. The RFA profile is a smoothed version of $v_0$ obtained by convolving along the graph structure; its variance is reduced by averaging over all paths. Let $v_0$ be a noisy indicator of peaks and $\mu_0 = \mathbb{E}[v_0]$. Then
\begin{equation}
\mathbb{E}[x_\gamma] = (I-\gamma P)^{-1} \mu_0, \qquad \mathrm{Var}[x_\gamma] = (I-\gamma P)^{-1} \mathrm{Var}[v_0] \big((I-\gamma P)^{-1}\big)^\top.
\end{equation}
Since $(I-\gamma P)^{-1}$ is a contraction-like operator on the subspace orthogonal to the stationary distribution, the variance of $x_\gamma$ is strictly smaller than that of $v_0$ in the relevant directions, while the mean is preserved. This is the graph analog of the bias-variance trade-off.

\textbf{Why does RFA anticorrelate with eQTL?} The smoothing argument above explains why RFA outperforms DC on promoter count, but it does \emph{not} predict anticorrelation with eQTL. Empirically, $\mathrm{AUC}(\mathrm{RFA}) = 0.386$ and $\mathrm{AUC}(-\mathrm{RFA}) = 0.614$ on eQTL, and $\mathrm{AUC}(\mathrm{RFA})$ decreases with eQTL strength (Section~5.3). We interpret this as follows: eQTL genes are expressed, but with stronger eQTL they are more \emph{buffered}, i.e., less dependent on 3D accessibility. RFA, being a smoothed accessibility score, is higher for genes whose expression is more strongly driven by 3D contacts, and lower for buffered genes with strong eQTL. This is a hypothesis, not a theorem, and it is stated as such in Section~6.3.

\textbf{Why does the combined model fail?} Let $X = \mathrm{RFA}$ and $Y = \mathrm{DC}$. Suppose $X = S(Y) + \varepsilon$ where $S$ is smoothing and $\varepsilon$ is noise. Then $\mathrm{AUC}(X+Y) \leq \mathrm{AUC}(X)$ whenever $\varepsilon$ dominates, because $Y$ reintroduces the noise $\varepsilon$ that $S$ removed. On eQTL, the combined model also fails ($\mathrm{AUC}=0.561$ vs $\mathrm{RFA}$ alone), but for a different reason: RFA and DC anticorrelate with eQTL in different ways, so $z(\mathrm{RFA}) + z(\mathrm{DC})$ partially cancels the anticorrelation.

\section{Data}

\subsection{Hi-C}
We use GM12878 in situ Hi-C (ENCODE/4DN, \texttt{4DNFIZGCA8AH.hic}) at 10\,kb resolution \cite{rao2014} for chromosomes 1, 17, and 19, extracted with \texttt{juicer\_tools dump observed NONE} and converted to dense \texttt{.npy} (int32) via a sparse-triplet converter. For MCF10A and MCF7, we use 40\,kb matrices for the same chromosomes (\texttt{MCF\{10A,7\}\_chr\{1,17,19\}\_40kb\_fixed.npy}). The 10\,kb matrices have 24\,895, 8\,325, and 5\,861 bins for chr1, chr17, and chr19 respectively, with nnz/$n^2$ of 0.012, 0.030, and 0.044; the sizes match chromosome lengths.

\subsection{TFBS (GM12878)}
Peaks are taken from \texttt{tfbs\_GM12878.bed} (BED6). We test 20 TFs: IKZF1, CTCF, SMC3, RAD21, PAX5, EBF1, EP300, POLR2A, YY1, GABPA, MAX, E2F4, ELF1, TCF12, BATF, BCL11A, IRF4, MEF2A, NRF1, and ZNF143.

\subsection{ATAC peaks (MCF10A / MCF7)}
For MCF7 we use IDR-thresholded peaks from GEO GSE169929 (\texttt{GSE169929\_ENCFF882OVP\_...bed.gz}); for MCF10A, broadPeak calls from GEO GSM3735772. Only the first three columns (chrom, start, end) are used; $v_0$ is the indicator of bins containing an ATAC peak. ATAC peaks mark accessible regulatory DNA \cite{thurman2012}.

\subsection{eQTL (GTEx v10 LCL)}
We use \texttt{QTD000221.all.tsv.gz} from the eQTL Catalogue (EBI), corresponding to GTEx v10 LCL (study \texttt{QTS000015}, dataset \texttt{QTD000221}), \texttt{quant\_method=ge}, \texttt{sample\_size=326}. The file is streamed in 2M-row chunks; for each gene we record $p_{\min}$ and $n_{\mathrm{sig}}$ (number of variants with $p<5\times10^{-8}$), and define $\mathrm{eqtl\_bin} = \mathbb{1}[n_{\mathrm{sig}}>0]$. The result is 18\,988 genes, of which 1\,567 (8.3\%) have at least one significant variant. The eQTL Catalogue provides uniformly reprocessed summary statistics
\cite{kerimov2021}.

\subsection{RNA (MCF10A / MCF7)}
We use \texttt{MCF10A\_RNA.tsv} (16\,829 genes) and \texttt{MCF7\_RNA.tsv} (19\,154 genes), with columns \texttt{gene\_id}, \texttt{gene\_name}, and \texttt{TPM}; there are no NAs and gene IDs are without version suffixes.

\subsection{Gene annotation}
Promoters are taken from \texttt{gene\_structure\_hg38.csv} (\texttt{element\_type=='promoter'}, first promoter per gene, \texttt{chrom} without \texttt{chr}). This yields 1043, 538, and 739 promoters on chr1, chr17, and chr19.

\subsection{Promoter TF count}
The promoter count matrix \texttt{matrix\_with\_hic\_enhancers.csv} contains 811 genes on chr1 with TF peak counts in the window $[-2000, +200]$\,bp from TSS (not score sums). It was built for chr1 only; for chr17/19 we merge directly with the eQTL target.

\section{Methods}

\subsection{Pipeline}
The pipeline is:
\begin{verbatim}
Hi-C .npy -> build_markov_matrix(topk=300, PC1-residualize) -> P
v0 (TF peaks / ATAC peaks) -> rfa_profile(P, v0, gamma=0.85) -> x_gamma
x_gamma[bin(gene)] -> RFA(gene)
\end{verbatim}

The same flow is shown schematically in Figure~\ref{fig:pipeline}.

\begin{figure}[htbp]
\centering
\includegraphics[width=\textwidth]{fig_pipeline.pdf}
\caption{Data-processing pipeline. Hi-C contacts are sparsified
(top-$k{=}300$), symmetrized and residualized against PC1 to build the
Markov transition matrix $P=D^{-1}A$. The RFA profile is the solution of
$(I-\gamma P)x=v_0$ with $\gamma=0.85$ and $v_0$ the TF/ATAC peak
indicator. Per-gene scores $\mathrm{RFA}(g)=x_{\gamma,b(g)}$ are tested
on eQTL (AUC of RFA and of $-$RFA) and on $\Delta$RNA (Spearman $\rho$,
MDR), against 1D distance, direct contact and PC1 baselines.}
\label{fig:pipeline}
\end{figure}

\subsection{Metrics}
We use ROC-AUC for binary targets; Spearman $\rho$ for continuous targets; partial Spearman by rank-residualizing both variables on controls; bootstrap CI (1000 resamples) for $\Delta$AUC and partial; permutation $p$ (500 permutations) for $\rho$ and AUC; and MDR (minimum detectable $\rho$) via Fisher $z$: $\tanh((z_\alpha + z_\beta)/\sqrt{n-3})$.

\textbf{Direction check.} Both $\mathrm{AUC}(\mathrm{RFA})$ and $\mathrm{AUC}(-\mathrm{RFA})$ are reported. If $\mathrm{AUC}(\mathrm{RFA})<0.5<\mathrm{AUC}(-\mathrm{RFA})$, RFA anticorrelates with the target. \textbf{No sign flip is applied} --- a flip would mask anticorrelation. In an earlier version of the pipeline, an automatic flip produced ``AUC=0.611''; in the final dose-response analysis the flip is removed.

\subsection{Controls}
We use 1D distance, direct contact (dc), PC1 (compartment), baseline RNA, permutation null, within-line sanity, and the direction check.

\section{Results}

\subsection{Promoter TF count (GM12878 chr1, 10\,kb, 811 genes)}
\begin{table}[h]
\centering
\begin{tabular}{lc}
\toprule
Metric & Value \\
\midrule
Mean AUC 1D & \textbf{0.976} \\
Mean AUC RFA & 0.836 \\
Mean AUC ChIPSP & 0.618 \\
Mean $\Delta$AUC(RFA $-$ 1D) & \textbf{$-0.140$} (0/20 TFs) \\
Mean partial RFA $\mid$ (1D, dc) & $+0.141$ (17/20 CI$>0$) \\
Mean partial RFA $\mid$ PC1 & $+0.348$ \\
\bottomrule
\end{tabular}
\end{table}

1D wins because the target (``is there a peak in the promoter $\pm2$\,kb'') and the 1D distance answer the same question. This is a methodological artifact, not a property of RFA. Partial $\mid$ 1D $= -0.068$: RFA does not add over 1D.

\subsection{eQTL: anticorrelation (GTEx v10 LCL, chr1+17+19, 10\,kb, 1625 genes, 124 positives)}
\begin{table}[h]
\centering
\begin{tabular}{lc}
\toprule
Metric & Value \\
\midrule
Mean AUC(RFA) & \textbf{0.386} (below chance) \\
Mean AUC($-$RFA) & \textbf{0.614} (above chance) \\
Mean AUC 1D & 0.422 \\
Mean AUC dc & 0.428 \\
$\rho_{\mathrm{RFA}}$ on $-\log_{10} p_{\min}$ & $\mathbf{-0.09}$ (19/20 negative) \\
partial RFA $\mid$ (1D, dc) & $-0.03$ (0/20 CI$>0$) \\
\bottomrule
\end{tabular}
\end{table}

RFA \emph{anticorrelates} with eQTL. High RFA corresponds to fewer eQTL. $\mathrm{AUC}(\mathrm{RFA})=0.386$ is below chance, which is the anticorrelation; $\mathrm{AUC}(-\mathrm{RFA})=0.614$ is above chance after sign inversion. No sign flip is applied. Genes with strong eQTL are less accessible in 3D (high RFA $\rightarrow$ fewer eQTL). A possible mechanism is \emph{expression buffering}: genes with strong eQTL have more buffered expression, less dependent on 3D accessibility. Alternative: RFA captures expression status (eQTL genes are expressed, but with weaker effects).

\subsection{eQTL dose-response: AUC(RFA) decreases with eQTL strength}
Figure~\ref{fig:dose} shows the mean dose-response curve; full tables for all 20 TFs are given in Appendix~S2. The mean $\mathrm{AUC}(Q4) - \mathrm{AUC}(Q1) = -0.056$, with 19/20 TFs showing a negative trend. RFA anticorrelates with eQTL strength: stronger eQTL corresponds to lower AUC(RFA). This is stable and confirms Section~5.2.

\begin{figure}[h]
\centering
\includegraphics[width=0.7\textwidth]{fig_dose.pdf}
\caption{Mean AUC(RFA) across eQTL strength quantiles (Q1 weak to Q4 strong). The line is the mean over 20 TFs; shaded area is $\pm$1 SD. AUC(RFA) decreases with eQTL strength (trend $-0.056$).}
\label{fig:dose}
\end{figure}

\subsection{$\Delta$RNA (MCF10A vs MCF7, 40\,kb, ATAC peaks, 1730 genes)}
\begin{table}[h]
\centering
\begin{tabular}{lc}
\toprule
Metric & Value \\
\midrule
$\rho(\Delta\mathrm{RFA}, \Delta\mathrm{RNA})$ & $\mathbf{+0.0047}$ \\
perm $p$ & 0.866 \\
null 95\% CI & $[-0.047, +0.050]$ \\
\textbf{MDR} ($n=1730$, power$=0.80$) & $\mathbf{0.067}$ \\
observed $|\rho|$ & 0.0047 (\textbf{14$\times$ below MDR}) \\
per-chr $\rho$ & chr1 $-0.031$, chr17 $+0.005$, chr19 $-0.000$ \\
$\rho(\mathrm{RFA}_{10A}, \mathrm{RNA}_{10A})$ & $-0.004$ ($p=0.87$) \\
$\rho(\mathrm{RFA}_{7}, \mathrm{RNA}_{7})$ & $+0.025$ ($p=0.30$) \\
partial $\mid$ baseline RNA & $-0.022$ \\
\bottomrule
\end{tabular}
\end{table}

$\Delta$RFA does not predict $\Delta$RNA. The effect is 14$\times$ below MDR --- absence of effect, not lack of power. Within-line sanity is also zero: RFA does not proxy expression at 40\,kb with ATAC peaks. Figure~\ref{fig:delta} shows the scatter.

\begin{figure}[h]
\centering
\includegraphics[width=0.6\textwidth]{fig_delta.pdf}
\caption{$\Delta$RFA vs $\Delta$RNA (MCF10A vs MCF7, 40\,kb, ATAC peaks, 1730 genes). The regression line is flat ($\rho=+0.005$), consistent with the absence of signal.}
\label{fig:delta}
\end{figure}

\section{Discussion}

\subsection{What RFA does}
RFA anticorrelates with eQTL at 10\,kb with TF peaks ($\mathrm{AUC}(-\mathrm{RFA})=0.614$, dose-response trend $-0.056$, 19/20 TFs). It is not a proxy for direct contact ($\rho(\mathrm{RFA},\mathrm{dc})=0.58$, 0/20 $>0.8$) and not a proxy for compartments (partial RFA $\mid$ PC1 $=+0.35$).

\subsection{What RFA does not do}
RFA does not predict $\Delta$RNA at 40\,kb with ATAC peaks ($\rho=0.005$, MDR$=0.067$). It does not add over 1D on promoter count (partial $\mid$ 1D $= -0.068$). It does not predict eQTL in the positive direction --- only anticorrelates.

\subsection{Why anticorrelation?}
Three non-exclusive hypotheses: (i) \textbf{expression buffering}: genes with strong eQTL have more stable expression, less dependent on 3D accessibility; (ii) \textbf{RFA captures expression status}: eQTL genes are expressed, and RFA correlates with expression (indirectly), but with eQTL strength inversely; (iii) \textbf{target artifact}: eQTL is a genetic layer, RFA is epigenetic, and the anticorrelation may reflect layer mismatch. We state all three explicitly, without choosing one without further data.

\subsection{Limitations}
First, resolution: 10\,kb is the working range; 40\,kb is not. Second, $v_0$: TF peaks give signal (anticorrelation); ATAC peaks do not. Third, target: eQTL is a genetic layer; $\Delta$RNA is dynamic and noisy. Fourth, PC1 residualization: post-hoc residualization does not remove PC1 (ratio$=0.623$). Fifth, single cell line/tissue: GM12878 + LCL; MCF10A/7. Sixth, sign flip: the automatic flip in the earlier pipeline masked anticorrelation; the final analysis uses explicit direction.

\subsection{Positioning}
``Continuous topological accessibility as a resolution-dependent 3D-aware GRN filter'': RFA anticorrelates with eQTL at 10\,kb with TF peaks and does not predict $\Delta$RNA at 40\,kb with ATAC. This is not a ``universal predictor'' but a filter with an inverted sign on eQTL and no signal on $\Delta$RNA. The applicability domain is static eQTL association at 10\,kb with TF peaks.

\section{Conclusion}
RFA anticorrelates with eQTL at 10\,kb with TF peaks ($\mathrm{AUC}(\mathrm{RFA})=0.386$, $\mathrm{AUC}(-\mathrm{RFA})=0.614$, dose-response trend $-0.056$). RFA does not predict $\Delta$RNA at 40\,kb with ATAC peaks ($\rho=0.005$, MDR$=0.067$). The anticorrelation may reflect expression buffering. The negative $\Delta$RNA result delineates the applicability boundary: RFA measures static accessibility, not dynamics. A methodological warning follows: a sign flip in AUC masks anticorrelation, so dose-response and explicit direction are mandatory.

\appendix
\section{Reproducibility}
Hi-C data were extracted with \texttt{juicer\_tools dump} and converted with \texttt{converter.py}. The $P$ matrix is built with \texttt{build\_markov\_matrix(topk=300, PC1-residualize)} and cached in \texttt{E:\textbackslash HIC\textbackslash results\textbackslash\{pilot\_multitf,eqtl\_multi\_chr,delta\_rfa\}}. The eQTL target is built by \texttt{eqtl\_target.py} (streaming GTEx v10 LCL). All scripts are listed in the companion code repository.

\section{Full dose-response tables}
Complete tables for all 20 TFs (4 quantiles of $-\log_{10} p_{\min}$ and 5 bins of $n_{\mathrm{sig}}$) are given in the supplementary CSV \texttt{eqtl\_dose\_response.csv}.

\section{Parameters}
$\gamma=0.85$, top-$k=300$, PC1 residualization, 10\,kb (GM12878) / 40\,kb (MCF). Bootstrap: 1000; permutation: 500. Minimum gene count per analysis: 20--50.

\section{Sign-flip correction}
The earlier pipeline applied \texttt{rfa\_for\_auc = -rfa if auc\_rfa\_neg > auc\_rfa else rfa}, yielding $\mathrm{AUC}=0.611$. Without the flip: $\mathrm{AUC}(\mathrm{RFA})=0.386$, $\mathrm{AUC}(-\mathrm{RFA})=0.614$. The final dose-response removes the flip. The article uses the no-flip values with explicit direction.


\begin{thebibliography}{9}

\bibitem{dekker2013}
Dekker J, Marti-Renom MA, Mirny LA.
Exploring the three-dimensional organization of genomes: interpreting chromatin interaction data.
\textit{Nat Rev Genet}. 2013;14(6):390--403.
doi:10.1038/nrg3454.

\bibitem{hu2013}
Hu M, Deng K, Qin Z, Liu JS.
Understanding spatial organizations of chromosomes via statistical analysis of Hi-C data.
\textit{Quant Biol}. 2013;1(2):156--174.
doi:10.1007/s40484-013-0016-0.

\bibitem{rao2014}
Rao SSP, Huntley MH, Durand NC, Stamenova EK, Bochkov ID, Robinson JT, et al.
A 3D map of the human genome at kilobase resolution reveals principles of chromatin looping.
\textit{Cell}. 2014;159(7):1665--1680.
doi:10.1016/j.cell.2014.11.021.

\bibitem{thurman2012}
Thurman RE, Rynes E, Humbert R, Vierstra J, Maurano MT, Haugen E, et al.
The accessible chromatin landscape of the human genome.
\textit{Nature}. 2012;489(7414):75--82.
doi:10.1038/nature11232.

\bibitem{kerimov2021}
Kerimov N, Hayhurst JD, Peikova K, Manning JR, Walter P, Kolberg L, et al.
A compendium of uniformly processed human gene expression and splicing quantitative trait loci.
\textit{Nat Genet}. 2021;53(9):1290--1299.
doi:10.1038/s41588-021-00924-w.

\end{thebibliography}

\end{document}